\subsection{非退化子模}

在本节及随后两节中，A 表示一个环，S 表示 A 中由不是 A 中零因子的元素组成的乘性子集，B 表示环 \(S^{-1}A\) ；A 典范地等同于 B 的一个子环（ \(\S\) 2, 第 1 节, 注 3）。因此 S 的元素在 B 中可逆。

对应用而言最重要的特殊情形之一是：A 是整环而 S 是 A 中由 \(\neq0\) 的元素组成的集合；这时 B 是 A 的分式域。

定义 3. 设 M 是 B 的一个 A-子模。若 B. M = B，则称 M 为非退化的。

若 B 是域，这个条件只是说 M 不化为 0。

命题 8. 设 M 是 B 的一个 A-子模。下列条件等价：

\begin{enumerate}
\def\labelenumi{(\alph{enumi})}
\item
  M 是非退化的。
\item
  M 与 S 相交。
\item
  若 \(j: \mathbf{M} \to \mathbf{B}\) 是典范单射，则同态 \(u = \mathbf{S}^{-1}j: \mathbf{S}^{-1}\mathbf{M} \to \mathbf{B}\) 是双射的。
  (a) 蕴含 (b)，因为若 B.M = B，则存在 \(a \in A\) 、\(s \in S\) 与 \(x \in M\) 使 \((a/s)x = 1\) ，于是 ax = s 属于 S ∩ M。为看出 (b) 蕴含 (c)，注意 u 已经是单射的（§2 第 4 节定理 1）；此外，若 \(x \in M \cap S\) ，则 \(x/x \in S^{-1}M\) 在 B 中 u 下的像等于 1，因此 u 是满射的。最后，(c) 显然蕴含 (a)。
\end{enumerate}

推论. 若 M 与 N 是 B 的两个非退化 A-子模，则 A-模 M + N、M.N 与 M ∩ N 都是非退化的。

结论对 \(\mathbf{M} + \mathbf{N}\) 是平凡的；另一方面，若 \(s \in \mathbf{S} \cap \mathbf{M}\) 且 \(t \in \mathbf{S} \cap \mathbf{N}\) ，则 \(st \in \mathbf{S} \cap (\mathbf{M}. \mathbf{N})\) 且 \(st \in \mathbf{S} \cap (\mathbf{M} \cap \mathbf{N})\) 。

给定 B 的两个 A-子模 M 与 N，用 N: M 表示由 \(b \in B\) 中满足 \(bM \subset N\) 的那些 b 组成的 B 的 A-子模（第 I 章 §2 第 10 节注）。若把每个 \(b \in N\) : M 对应到 M 到 N 的同态 \(h_{b}: x \mapsto bx\) ，就得到典范同态 \(b \mapsto h_{b}\) ：N: M → \(\operatorname{Hom}_{\mathbf{A}}(\mathbf{M}, \mathbf{N})\) 。

命题 9. 设 M、N 是 B 的两个 A-子模。若 M 是非退化的，则从 N: M 到 \(\mathrm{Hom}_{\mathbf{A}}(\mathbf{M},\mathbf{N})\) 的典范同态是双射的。

设 \(s \in S \cap M\) 。若 \(b \in N: M\) 对所有 \(x \in M\) 满足 bx = 0，则 bs = 0，由于 s 在 B 中不是零因子，得 b = 0。另一方面，设 \(f \in \operatorname{Hom}_{\mathbf{A}}(\mathbf{M}, \mathbf{N})\) ，令 \(b = f(s) / s\) ；对每个 \(x \in \mathbf{M}\) ，存在 \(t \in \mathbf{S}\) 使 \(tx \in \mathbf{A}\) 。于是

\[
f (x) = s ^ {- 1} t ^ {- 1} f (s t x) = s ^ {- 1} t ^ {- 1} t x f (s) = b x,
\]

由此 \(b \in \mathbf{N}\) : M 且 \(f = h_b\) ，命题得证。

注. 特别地，A: M 典范地等同于 M 的对偶模 M*，M × M* 上的典范双线性型等同于乘法 B × B → B 在 M × (A: M) 上的限制。

